We use the projective change of variables from \cite{r-75}, Theorem 4.10, and \cite{r-30}, Section 2.2. A finite set of rotations reduces the direction supremum to \(|\omega_n|\ge c_n\) for some \(c_n>0\). Translate any unit cube \(Q\) to the origin and write \(x=(x',t_0)\in Q\). For \(\rho\) sufficiently small, each tube in the restricted supremum lies in \(5Q\), in a slab \(c_n/4\le |s-t_0|\le2\) with a fixed sign of \(s-t_0\). Fix the height \(t_0\) and one sign. On this bounded slab and its image the map
\[
 (y',s)\longmapsto (Y',Y_n)=\left(y'/(s-t_0),\,1/(s-t_0)\right)
\]
and its inverse are uniformly Lipschitz, with Jacobian determinants bounded above and below. On the axis \(y'=x'+(s-t_0)v\) of any such tube the transformed line is \(Y'=v+Y_nx'\). Its direction is \((x',1)/|(x',1)|\) independently of \(v\). The image of the axis segment has bounded length and the image of the tube lies in its \(C_n\rho\)-neighborhood. Hence this image is covered by a bounded number of unit tubes of radius \(C_n\rho\) in that direction. The volumes of the original tube and of each covering tube are comparable to \(\rho^{n-1}\).

Transport \(|f|\) restricted to \(5Q\) and the fixed sign slab by the projective map and extend by zero. The resulting \(F\) depends on \(t_0\) and sign and works for all \(x'\) in the slice. Its \(L^n\) norm is \(\lesssim\|f\|_{L^n(5Q)}\) by the Jacobian bounds. Each original average is bounded by a constant times \(K_{C_n\rho}F((x',1)/|(x',1)|)\), using the measure comparison. The direction map is injective and has area Jacobian bounded below for the bounded \(x'\) in this cube. Apply the assumed Kakeya estimate to bound the slice \(L^n\) norm by \(C L(1/\rho)^{2A}\|f\|_{L^n(5Q)}\). Sum over the two signs and integrate in \(t_0\). Then translate over unit lattice cubes and use bounded overlap of their dilates \(5Q\). For \(\rho\) bounded below, the tubes lie in a fixed ball about \(x\) and have volume bounded below, so Young's inequality gives the asserted estimate directly.
